% GENERATED FILE. Edit canonical lessons, then run npm run generate:books.
% canonical-source-hash: fnv1a64:5da456724ffe6dd5
% canonical-lessons: LA-C526-sano, LA-C526-medeor, LA-C526-convalesco, LA-C526-recreo, LA-C526-potio

\chapter{Healing}
\label{ch:la-healing}
\hlchaptermodality{526}

\begin{chapteropening}
\textbf{By the end of this chapter:} \emph{I can talk about healing and getting well.}

Complete the last lesson of chapter 526: I can talk about healing and getting well.
\end{chapteropening}

\section[sānō, sānāre]{sānō, sānāre — to heal, cure}

\label{lesson:LA-C526-sano}



\begin{quote}
Before the new one: say the Latin for bleary-eyed, then the Latin for inflammation of the eyes.
\end{quote}

\subsection*{You'll want to know: sānō, sānāre}
\textbf{sānō, sānāre} — \textquotedblleft{}to heal, cure\textquotedblright{}.

\begin{quote}
Medicus vulnus mīlitis sānat.
\end{quote}

\begin{quote}
\emph{The doctor heals the soldier's wound.}
\end{quote}

\emph{sānō, sānāre, sānāvī, sānātum}, from \emph{sānus}. English \emph{sanatorium}.

\begin{cousinweb}
\textbf{sānō} is built on \emph{sānus} \textquotedblleft{}sound, healthy\textquotedblright{} (Lewis \& Short; EtymDB). English \emph{sanatorium} comes from it.
\end{cousinweb}

\begin{grammarlens}[title={What you've built}]
One more word for telling, explaining and describing at B1.
\end{grammarlens}

\subsection*{Your turn}
\begin{itemize}
\raggedright
  \item \emph{Say it:} \emph{sānō, sānāre}
  \item \emph{Say it:} \emph{sānō, sānāre}, once more
  \item \emph{Say it:} a sentence of your own with \emph{sānō, sānāre}
  \item \emph{Recall:} say the Latin for bleary-eyed, then the Latin for inflammation of the eyes, then say \emph{sānō, sānāre} again
\end{itemize}

\begin{tcolorbox}[breakable,colback=teal!4,colframe=teal!35!black,title={Before you move on}]
What does sānō, sānāre mean? (\textquotedblleft{}To heal, cure\textquotedblright{}.) Say it once more.
\end{tcolorbox}
\section[medeor, medērī]{medeor, medērī — to remedy, be good for}

\label{lesson:LA-C526-medeor}



\begin{quote}
Before the new one: say the Latin for inflammation of the eyes, then the Latin for to heal.
\end{quote}

\subsection*{You'll want to know: medeor, medērī}
\textbf{medeor, medērī} — \textquotedblleft{}to remedy, be good for\textquotedblright{}.

\begin{quote}
Haec herba morbō medētur.
\end{quote}

\begin{quote}
\emph{This herb is good for the disease.}
\end{quote}

Deponent, with the dative: \emph{morbō medērī}. No perfect in use. \emph{Medicus}, \emph{medicīna} and \emph{remedium} come from the same root.

\begin{cousinweb}
\textbf{medeor} is inherited through Proto-Italic from a Proto-Indo-European root \emph{med-} \textquotedblleft{}to measure, take care\textquotedblright{} (EtymDB). \emph{medicus}, \emph{medicīna} and \emph{remedium} belong to it, and so English \emph{medicine} and \emph{remedy}.
\end{cousinweb}

\begin{grammarlens}[title={What you've built}]
One more word for telling, explaining and describing at B1.
\end{grammarlens}

\subsection*{Your turn}
\begin{itemize}
\raggedright
  \item \emph{Say it:} \emph{medeor, medērī}
  \item \emph{Say it:} \emph{medeor, medērī}, once more
  \item \emph{Say it:} a sentence of your own with \emph{medeor, medērī}
  \item \emph{Recall:} say the Latin for inflammation of the eyes, then the Latin for to heal, then say \emph{medeor, medērī} again
\end{itemize}

\begin{tcolorbox}[breakable,colback=teal!4,colframe=teal!35!black,title={Before you move on}]
What does medeor, medērī mean? (\textquotedblleft{}To remedy, be good for\textquotedblright{}.) Say it once more.
\end{tcolorbox}
\section[convalēscō, convalēscere]{convalēscō, convalēscere — to recover, get well}

\label{lesson:LA-C526-convalesco}



\begin{quote}
Before the new one: say the Latin for to heal, then the Latin for to remedy.
\end{quote}

\subsection*{You'll want to know: convalēscō, convalēscere}
\textbf{convalēscō, convalēscere} — \textquotedblleft{}to recover, get well\textquotedblright{}.

\begin{quote}
Post longum morbum avus convalēscit.
\end{quote}

\begin{quote}
\emph{After a long illness grandfather is getting well.}
\end{quote}

\emph{convalēscō, convalēscere, convaluī}, built on \emph{valeō} — to be well. English \emph{convalesce}.

\begin{cousinweb}
\textbf{convalēscō} is built on \emph{con-} and \emph{valēscō}, the inceptive of \emph{valeō} \textquotedblleft{}to be well, strong\textquotedblright{} (Lewis \& Short). English \emph{convalesce} comes from it.
\end{cousinweb}

\begin{grammarlens}[title={What you've built}]
One more word for telling, explaining and describing at B1.
\end{grammarlens}

\subsection*{Your turn}
\begin{itemize}
\raggedright
  \item \emph{Say it:} \emph{convalēscō, convalēscere}
  \item \emph{Say it:} \emph{convalēscō, convalēscere}, once more
  \item \emph{Say it:} a sentence of your own with \emph{convalēscō, convalēscere}
  \item \emph{Recall:} say the Latin for to heal, then the Latin for to remedy, then say \emph{convalēscō, convalēscere} again
\end{itemize}

\begin{tcolorbox}[breakable,colback=teal!4,colframe=teal!35!black,title={Before you move on}]
What does convalēscō, convalēscere mean? (\textquotedblleft{}To recover, get well\textquotedblright{}.) Say it once more.
\end{tcolorbox}
\section[recreō, recreāre]{recreō, recreāre — to refresh, restore}

\label{lesson:LA-C526-recreo}



\begin{quote}
Before the new one: say the Latin for to remedy, then the Latin for to recover.
\end{quote}

\subsection*{You'll want to know: recreō, recreāre}
\textbf{recreō, recreāre} — \textquotedblleft{}to refresh, restore\textquotedblright{}.

\begin{quote}
Somnus corpus fessum recreat.
\end{quote}

\begin{quote}
\emph{Sleep refreshes the tired body.}
\end{quote}

\emph{recreō, recreāre, recreāvī, recreātum}: \emph{re-} 'again' + \emph{creō} 'to make'. English \emph{recreation}.

\begin{cousinweb}
\textbf{recreō} is built on \emph{re-} \textquotedblleft{}again\textquotedblright{} and \emph{creō} \textquotedblleft{}to create\textquotedblright{} (Lewis \& Short; EtymDB): to make new again. English \emph{recreation} comes from it.
\end{cousinweb}

\begin{grammarlens}[title={What you've built}]
One more word for telling, explaining and describing at B1.
\end{grammarlens}

\subsection*{Your turn}
\begin{itemize}
\raggedright
  \item \emph{Say it:} \emph{recreō, recreāre}
  \item \emph{Say it:} \emph{recreō, recreāre}, once more
  \item \emph{Say it:} a sentence of your own with \emph{recreō, recreāre}
  \item \emph{Recall:} say the Latin for to remedy, then the Latin for to recover, then say \emph{recreō, recreāre} again
\end{itemize}

\begin{tcolorbox}[breakable,colback=teal!4,colframe=teal!35!black,title={Before you move on}]
What does recreō, recreāre mean? (\textquotedblleft{}To refresh, restore\textquotedblright{}.) Say it once more.
\end{tcolorbox}
\section[pōtiō, pōtiōnis]{pōtiō, pōtiōnis — a draught, drink; a potion}

\label{lesson:LA-C526-potio}



\begin{quote}
Before the new one: say the Latin for to recover, then the Latin for to refresh.
\end{quote}

\subsection*{You'll want to know: pōtiō, pōtiōnis}
\textbf{pōtiō, pōtiōnis} — \textquotedblleft{}a draught, drink; a potion\textquotedblright{}.

\begin{quote}
Medicus aegrō pōtiōnem amāram dat.
\end{quote}

\begin{quote}
\emph{The doctor gives the patient a bitter draught.}
\end{quote}

Feminine, third declension, from the root of \emph{pōtō} — to drink. English \emph{potion}, and \emph{poison} too.

\begin{cousinweb}
\textbf{pōtiō} is built on the root of \emph{pōtō} \textquotedblleft{}to drink\textquotedblright{} (Lewis \& Short; EtymDB). English \emph{potion} comes from it, and \emph{poison} too, through Old French.
\end{cousinweb}

\begin{grammarlens}[title={What you've built}]
One more word for telling, explaining and describing at B1.
\end{grammarlens}

\subsection*{Your turn}
\begin{itemize}
\raggedright
  \item \emph{Say it:} \emph{pōtiō, pōtiōnis}
  \item \emph{Say it:} \emph{pōtiō, pōtiōnis}, once more
  \item \emph{Say it:} a sentence of your own with \emph{pōtiō, pōtiōnis}
  \item \emph{Recall:} say the Latin for to recover, then the Latin for to refresh, then say \emph{pōtiō, pōtiōnis} again
\end{itemize}

\begin{tcolorbox}[breakable,colback=teal!4,colframe=teal!35!black,title={Before you move on}]
What does pōtiō, pōtiōnis mean? (\textquotedblleft{}A draught, drink; a potion\textquotedblright{}.) Say it once more.
\end{tcolorbox}
